% scope_and_limitations.tex
% Insertable section for eci.tex
% Author: Kevin Remondière (drafted by Claude Sonnet 4.6 agent, 2026-05-03)
% Compile: standalone or % scope_and_limitations.tex
% Insertable section for eci.tex
% Author: Kevin Remondière (drafted by Claude Sonnet 4.6 agent, 2026-05-03)
% Compile: standalone or % scope_and_limitations.tex
% Insertable section for eci.tex
% Author: Kevin Remondière (drafted by Claude Sonnet 4.6 agent, 2026-05-03)
% Compile: standalone or % scope_and_limitations.tex
% Insertable section for eci.tex
% Author: Kevin Remondière (drafted by Claude Sonnet 4.6 agent, 2026-05-03)
% Compile: standalone or \input{scope_and_limitations} after last body section.
%
% This section documents the bounded scope of the ECI framework.
% It is NOT a defeat — it is a calibration statement.
% Every paragraph is written to pass the critic's smell test:
%   "Can a hostile referee accuse this of overclaiming?"
% If the answer is yes, the paragraph has been rewritten.

\section{Scope and Limitations of the ECI Framework}
\label{sec:scope}

The Entanglement--Complexity--Information (ECI) programme is a
mathematical-physics framework in the post-CLPW algebraic-quantum-gravity
tradition~\cite{Chandrasekaran2023,Chandrasekaran2022}.
It constructs a type-II$_\infty$ von Neumann algebra for FRW cosmological
diamonds via a quantum-reference-frame crossed product, derives a modular
(Krylov) complexity measure on that algebra, and uses the Bisognano--Wichmann
theorem and Hadamard-state theory to characterize which boundary conditions
are algebraically admissible near the Big Bang.  The framework's value is
precision and exclusion: it gives sharp structural results about
\emph{what is and is not algebraically possible} in cosmological quantum
field theory on a fixed background.

This section states, with equal emphasis, what ECI claims and what it does
not claim.  The negative results listed in \S\ref{subsec:nogo} are presented
as successes of calibration, not as failures of the programme.  A framework
that has explicitly tested and retired seventeen candidate bridges to
phenomenology has narrowed its own scope in a scientifically productive way.

% ---------------------------------------------------------------------------
\subsection{What ECI Claims and Is Positioned to Prove}
\label{subsec:claims}
% ---------------------------------------------------------------------------

\textbf{Core algebraic apparatus.}  The following results are established
within the framework, with the level of proof noted explicitly.

\begin{enumerate}

\item \textbf{FRW type-II$_\infty$ crossed product (Theorems 3.5--3.6,
  \texttt{paper/frw\_typeII\_note}).}
  The local observable algebra on a comoving past-light-cone diamond in
  radiation-era or matter-era FRW is a type-II$_\infty$ von Neumann factor
  when constructed as a quantum-reference-frame crossed product $\mathcal{A}(D_O)
  \rtimes_\sigma \mathbb{R}$.  The construction follows Witten~\cite{Witten2022},
  Chandrasekaran--Penington--Witten~\cite{Chandrasekaran2023}, and
  Leutheusser--Liu~\cite{Leutheusser2022}; the FRW adaptation is due to this
  work.  The factor carries a canonical tracial weight, but this trace fixes
  only a \emph{relative} scale (not an absolute energy density).

\item \textbf{Conformal pullback of the stress tensor (Lemma 3.3,
  sympy-verified).}
  The unitary $U$ implementing the FRW$\to$Minkowski conformal rescaling
  $a(\eta)$ satisfies $U^{-1}\,T^{\rm Mink}_{00}\,U = a^2(\eta)\,
  \mathcal{T}^{\rm FRW}_{00}$ off-shell, to all orders in the Taylor-mode
  expansion of the test-function support.  The conjugation is the
  backbone of the $K_{\rm FRW}$ closed-form derivation.

\item \textbf{$K_{\rm FRW}$ closed form (R-piste 6, $\sim$80\% paper-ready).}
  For a massless conformally-coupled scalar on a radiation-era FRW
  past-light-cone diamond $D_R$ of comoving radius $R_{\rm conf}$,
  the modular Hamiltonian is
  \begin{equation}
    K_{\rm FRW} \;=\; \frac{\pi}{R_{\rm conf}}
    \int_{D_R} d^3x\;
    \bigl[R_{\rm conf}^2 - (\eta - \eta_c)^2 - |\mathbf{x}-\mathbf{x}_c|^2\bigr]\,
    a^2(\eta)\,\mathcal{T}^{\rm FRW}_{00},
    \label{eq:kfrw}
  \end{equation}
  derived by conformal pullback from the Hislop--Longo
  Minkowski-diamond formula~\cite{HislopLongo1982}.  The prefactor
  $\pi/R_{\rm conf}$ is structural (Bisognano--Wichmann unimodular
  normalization $\beta=2\pi$, rescaled by the CHM08 Unruh relation); it is
  \emph{not} matched post-hoc to any observational value.

\item \textbf{Krylov-Diameter correspondence (Theorem 4, conditional).}
  In the late-modular-time asymptotic regime, the Krylov spread complexity
  $C_k(t)$ on the type-II$_\infty$ crossed-product algebra satisfies
  \begin{equation}
    \frac{1}{C_k(t)}\frac{dC_k}{dt} \;=\; \frac{1}{R_{\rm proper}(\eta_c)}
    + O(C_k^{-1}),
    \label{eq:krylov_rate}
  \end{equation}
  where $R_{\rm proper}(\eta_c) = a(\eta_c)\,R_{\rm conf}$ is the gauge-invariant
  proper radius of the diamond at conformal time $\eta_c$.  This identity is
  gauge-invariant (both sides are spacetime scalars) and is verified
  symbolically era-by-era (sympy).
  \textbf{Caveat:} the result is conditional on the \emph{restricted-class}
  Lanczos asymptotic $b_n = \pi n$ (Block A1, Open Question~1 in
  \texttt{krylov\_diameter.tex}) holding for the modular Hamiltonian's
  spectral class on a type-II$_\infty$ KMS state.  This transfer from the
  type-III$_1$ / CFT$_2$ setting (where the universal operator-growth
  hypothesis is established) to type-II$_\infty$ is \emph{not yet proved}.
  The conditional result is clean and publishable; the unconditional result
  requires Block A1 closure.

\item \textbf{Algebraic Weyl Curvature Hypothesis — Big-Bang
  non-Hadamard boundary (T2-FRW and T2-Bianchi I, $\sim$80\%
  paper-ready).}
  No state in the BFV--Verch Hadamard folium of $\Omega_{\rm FRW}$ extends
  to $\eta = 0$ in flat FRW (Theorem T2-FRW, three independent
  sympy-verified obstructions) or in Bianchi I (Theorem T2-Bianchi I, two
  independent IR divergences).  The physics consequence (Theorem T3) is that
  the algebras $({\mathcal A}(D_\infty), \Omega_{\rm FRW})$ and
  $(\mathcal{A}(D_{\rm BB}), \Omega_{\rm FRW})$ are not isomorphic, giving a
  rigorous algebraic correlate of the Penrose--Weyl curvature hypothesis.
  \textbf{Convention caveat:} the result is stated within the Hollands--Wald
  Hadamard microlocal convention~\cite{HollandsWald2001,BFV2003};
  alternative UV regulators may alter the IR divergence structure.
  The result does \emph{not} constitute a derivation of the
  cosmological arrow of time from first principles: the low-complexity
  initial condition (Past Hypothesis) is assumed as a boundary condition,
  not derived.

\item \textbf{Levier 1B MCMC (observational, not algebraic).}
  A joint MCMC run over Planck 2018 + ACT DR6 + DESI DR2 + Pantheon+ +
  KiDS-1000, with non-minimal coupling $\xi_\chi$ free, returns:
  $H_0 = 67.67 \pm 0.57$ km/s/Mpc, $\xi_\chi = -0.00003$
  $[-0.0164, +0.0163]$ (68\% CL), $\log B_{10} = -1.37$ (Savage--Dickey,
  weak/moderate evidence for $\Lambda$CDM over NMC).
  The ECI programme provided the NMC prior structure; the MCMC is an
  observational constraint, not an algebraic theorem.

\end{enumerate}

% ---------------------------------------------------------------------------
\subsection{What ECI Explicitly Does NOT Claim: the NO-GO Catalogue}
\label{subsec:nogo}
% ---------------------------------------------------------------------------

The following seventeen candidate connections between the ECI apparatus and
physical observables have been tested and returned negative or null results.
Each entry is a success of calibration.  They are listed here explicitly so
that future readers, referees, and automated agents do not re-open avenues
that have already been closed with stated reasons.

\subsubsection*{The five Pistes (systematic programme failures)}

\begin{description}

\item[Piste 1 (partial positive only).]
  $H = \text{Krylov rate}$: the identity
  $\tfrac{1}{C_k}\tfrac{dC_k}{dt} = H(t)$ holds exactly for a
  radiation-era Hubble-horizon diamond (gauge-invariant, sympy-verified) and
  is off by a factor of $1/2$ for a matter-era past-light-cone diamond.
  The match in the de Sitter and Hubble-horizon cases is tautological
  (the diamond is defined to have $R_{\rm proper} = H^{-1}$).
  \emph{ECI does not claim a general derivation of the Hubble parameter
  from Krylov complexity.}

\item[Piste 2 (NEGATIVE).]
  Era transitions as Jones--Longo subfactor inclusions: the algebra
  inclusion at the radiation-to-matter transition has Jones index equal to 1
  (the factors are isomorphic) and the subfactor is trivial.
  \emph{No cosmological transition amplitude is derived from subfactor theory.}

\item[Piste 3 (TAUTOLOGICAL).]
  $\Lambda$ from modular relaxation: the type-II$_\infty$ tracial weight
  fixes only a relative scale; equating it to $\rho_\Lambda / M_P^4 \approx
  10^{-121}$ requires a 121-order normalisation choice that is not derivable
  from the algebra alone.  The trace is gauge-fixed by an arbitrary modular
  normalisation, not by a physical energy density.
  \emph{ECI does not contribute a derivation of the cosmological constant.}
  (See also \S\ref{subsec:outside}.)

\item[Piste 4 (NO-GO).]
  Riemann zeros as cosmological modular Hamiltonian spectrum: the spectrum
  of $K_{\rm FRW}$ (eq.~\eqref{eq:kfrw}) is continuous, not discrete;
  the Riemann zeros form a discrete set.  A continuous-to-discrete
  mapping would require compactification not present in the apparatus.
  \emph{ECI does not offer a route to the Riemann Hypothesis via cosmology.}

\item[Piste 5 (NO-GO).]
  SM field content from modular 2-categorical compatibility: the
  type-I (finite-dimensional) algebra arising from any fixed Fock-space
  sector has only \emph{inner} modular automorphisms; the SM gauge group
  cannot be read off from the modular automorphism group of a type-I factor.
  \emph{ECI does not derive the Standard Model field content.}

\end{description}

\subsubsection*{The ten R-pistes (targeted rapid explorations)}

\begin{description}

\item[R1 (NUMEROLOGY).]
  $\Lambda$ from T1$\leftrightarrow$T2 algebraic asymmetry: the numerical
  match relied on a post-hoc normalisation; no structural reason for the
  asymmetry to equal $\rho_\Lambda$.

\item[R2 (POSITIVE BYPRODUCT, not $H_0$).]
  $H_0$ from Krylov-Diameter Theorem 4, Block A: the sympy-verified
  $\langle p_3 \rangle$ closed form is a genuine result, but it does not
  constitute a derivation of $H_0$.  The Krylov rate equals $1/R_{\rm proper}$,
  not $H_0$, in the matter era (factor of $1/2$ discrepancy;
  \S\ref{subsec:claims} item 4).

\item[R3 (KILLED: Verch 1994).]
  $S_8$ as BFV folium-projection ambiguity: Verch 1994~\cite{Verch1994}
  proves local quasi-equivalence of all Hadamard states; the alleged
  ``projection ambiguity'' does not exist within the folium.
  Levier 1B MCMC confirms null: NMC alone does not reduce the
  $S_8$ tension at Planck + KiDS-1000 sensitivity.

\item[R4 (REFUTED: Lie algebra).]
  SM three generations from Bianchi IX $SU(2)$ Kasner axes: the
  dimension count of the $\mathfrak{su}(2)$ BKL algebra is 3 (the three
  Kasner directions), but this is the rank of the spatial isometry group,
  not a fermion-generation count.  The Lie-algebraic bridge is not present.

\item[R5 (CONDITIONAL COMPANION, downgraded).]
  Riemann Hypothesis via Maass--Selberg + theta-correspondence:
  the earlier framing as an RH-proof has been retracted.
  The surviving claim is a \emph{conditional companion}: if the
  BFV-folium $\leftrightarrow$ CCM Eisenstein spectral correspondence holds,
  certain spectral conditions on the modular Hamiltonian would follow.
  No unconditional RH implication is claimed.

\item[R6 (CONDITIONAL COMPANION).]
  AdS/CFT in FRW type-II$_\infty$ frame: the original framing
  (``ECI establishes AdS/CFT for FRW'') is falsified.  The surviving
  claim is structural: the FRW type-II$_\infty$ factor is compatible
  with CLPW-style observer algebras, but the full holographic dictionary
  in an FRW background is not established.

\item[R7 (NUMEROLOGY).]
  Hierarchy from $c'$ modular trace ratios: the numerical match between
  the $c'$ parameter and $M_W/M_P$ depends on the same arbitrary
  normalisation identified in Piste 3.  No structural derivation.

\item[R8 (REFUTED: Weyl law).]
  Generation count from PSL$(2,\mathbb{Z})\backslash\mathbb{H}$ Eisenstein/cusp
  sectors: the Weyl law on the modular surface gives a
  \emph{continuous} eigenvalue density $\sim T \log T$, not a finite
  count of three.  The sector-counting route is blocked.

\item[R9 (KILLED: Theorem T2).]
  Dark matter from Bianchi VIII Mixmaster Kasner-bounce relics:
  Theorem T2-Bianchi I (and its extension to Bianchi VIII type-A)
  shows that no Hadamard state extends to the initial Kasner singularity;
  the bounce-relic mechanism requires initial conditions in the folium,
  which do not exist.

\item[R10 (NO-GO: DKN walls).]
  SM gauge from BKL $\leftrightarrow$ $E_{10}$ Kac--Moody:
  the Damour--Kleinschmidt--Nicolai billiard walls encode the
  $E_{10}$ structure but not the SM gauge factors $SU(3)\times SU(2)
  \times U(1)$; the coset decomposition does not produce SM multiplets.

\end{description}

% ---------------------------------------------------------------------------
\subsection{Open Problems Within ECI's Scope}
\label{subsec:open_within}
% ---------------------------------------------------------------------------

The following problems are within the ECI apparatus's reach and represent
the programme's live research frontier.  They are listed with their
current readiness level.

\begin{enumerate}

\item \textbf{Block A1 — unconditional UOGH transfer (open).}
  Proving $b_n = \pi n$ as an asymptotic on a type-II$_\infty$ KMS state
  without appealing to the universal operator-growth hypothesis in the
  type-III$_1$ / CFT$_2$ setting.  Completion would remove the ``conditional''
  caveat from Theorem 4.
  \emph{Status:} restricted-class Vanlessen/Levin--Lubinsky route in progress
  (B6 wave-3); full closure estimated at 2--3 months.

\item \textbf{Connes-embedding status of the FRW II$_\infty$ factor (open).}
  Is the crossed-product factor $\mathcal{A}(D_O) \rtimes_\sigma \mathbb{R}$
  approximately finite-dimensional (hyperfinite)?
  By Connes--Takesaki 1977 the answer is almost certainly yes (a crossed
  product of a hyperfinite type-III$_1$ factor by $\mathbb{R}$ is hyperfinite
  II$_\infty$), which would imply every cosmological observer in the BFV
  folium sees the same unique factor.
  \emph{Status:} 6--8 week short note; math.OA submission.

\item \textbf{T2-Bianchi V and T2-Bianchi IX extensions (open).}
  Theorem T2-Bianchi I is proved.  Extension to Bianchi V (negatively-curved
  spatial sections) is conditional on Hadamard-state existence on $H^3$;
  extension to Bianchi IX requires the BKL invariant-measure analysis of the
  Mixmaster attractor.
  \emph{Status:} partial proofs in \texttt{notes/T2\_bianchi\_2026\_05\_03\_jour2};
  3--8 weeks per case.

\item \textbf{BEC analogue-FRW experimental falsifier (open).}
  Lemma 3.3 (conformal pullback) predicts a modular spectral density on a
  BEC analogue of the FRW past-light-cone diamond.  Comparison against
  BEC entanglement-Hamiltonian tomography (Dalmonte--Vermersch--Zoller 2018
  framework) would constitute the first ECI experimental falsifier outside
  cosmological MCMC.
  \emph{Status:} 6--9 months; requires a BEC-experimental co-author.

\item \textbf{Algebraic WCH paper submission (near-ready).}
  The \texttt{paper/algebraic\_wch\_bianchi/note.tex} is 80\% paper-ready
  for Foundations of Physics.  Theorem T2-Bianchi I is the central result;
  the paper requires 2--3 weeks of polish.

\item \textbf{Bianconi 2025 PRD Comment (near-ready).}
  A 4-page note pointing out a Lorentzian-signature gap in Bianconi's
  entropic-gravity framework~\cite{Bianconi2025} has been drafted.
  The comment does \emph{not} claim ECI resolves dark energy or $\Lambda$;
  it is a structural mathematical objection
  ($\mathrm{Tr}(\tilde{g}\ln\tilde{g}^{-1}) \neq 0$ in the Lorentzian case)
  and an identification of where the ECI algebraic approach diverges from
  the entropic approach.
  \emph{Status:} 1--2 weeks to submission.

\end{enumerate}

% ---------------------------------------------------------------------------
\subsection{Open Problems Explicitly Outside ECI's Apparatus Reach}
\label{subsec:outside}
% ---------------------------------------------------------------------------

The following problems are \emph{outside} the ECI apparatus's reach.
This is not a temporary limitation to be overcome by adding more structure:
in each case a structural reason is given for why the apparatus does not
and cannot couple.

\begin{description}

\item[Cosmological constant $\Lambda \approx 10^{-122}\,M_P^4$.]
  The type-II$_\infty$ tracial weight fixes a relative scale, not an absolute
  energy density.  Equating the trace to $\rho_\Lambda / M_P^4$ introduces a
  normalisation that must be fixed from outside the algebra.
  Three independent arithmetic checks (NCG spectral cutoff, $E_8$ instanton,
  ECI internal $\kappa_R / C_k^{\rm max}$) each fail by 14--89 orders of
  magnitude~\cite{ECI_cosmo_const_audit}.
  \emph{No ECI paper claims a derivation of $\Lambda$.}

\item[Hubble tension ($H_0 = 73.04 \pm 1.04$ vs $67.67 \pm 0.57$
  km/s/Mpc, 4.5$\sigma$, UNRESOLVED).]
  Levier 1B (NMC $\xi_\chi$ free, full Planck + ACT + DESI + Pantheon+
  + KiDS-1000) returns $\log B_{10} = -1.37$ (weak/moderate evidence
  \emph{for} $\Lambda$CDM over NMC).  The NMC sector does not move $H_0$.
  Algebraically, $K_{\rm FRW}$ contains no Hubble-flow normalisation;
  Theorem 4 gives $1/R_{\rm proper}$, not $H_0$, as the Krylov rate in the
  matter era.
  \emph{ECI does not claim to resolve the Hubble tension.
  The 4.5$\sigma$ discrepancy remains open.}

\item[Standard Model particle content and gauge structure.]
  Piste 5 (inner modular automorphisms) and R4, R8, R10 (various
  generation-counting routes) all return NO-GO.  The modular automorphism
  group of any type-I sector is inner; the SM gauge factors are not
  extractable from the BFV-folium algebra.
  \emph{ECI does not predict or derive the SM gauge group, particle masses,
  or fermion generation count.}

\item[$S_8$ tension (KiDS vs Planck, $\sim$2.6$\sigma$).]
  Wolf + 2025 NMC mechanism and Levier 1B both return null on $S_8$
  (R-piste S$_8$ NEGATIVE).  Adding $\sigma_8$ as a derived parameter in
  a future MCMC run is advised before closing this entry, but NMC alone
  is not expected to resolve the tension.

\item[Swampland conjectures and string-landscape Λ-selection.]
  The Swampland $\times$ NMC cross-constraint (D8) tightens $|\xi|$ by
  $\approx 250\times$ and shows that the ECI $\chi$ field cannot also be the
  dark-dimension bulk mode (A4 $\neq$ A5).  Beyond this \emph{constraint},
  ECI does not contribute to swampland physics, de Sitter stability, or
  landscape statistics.

\item[Galaxy-scale structure (JWST early massive galaxies, $z > 10$).]
  ECI's apparatus operates on homogeneous--isotropic local diamonds.
  It has no coupling to halo abundance, star-formation efficiency, or
  first-light timing.  No ECI result bears on the JWST stellar-mass excess.

\item[SM precision observables (W-mass, muon $g-2$, lepton universality).]
  ECI's operator-algebra machinery has no coupling to loop corrections in
  perturbative QED/QCD.  The catalogue records no-apparatus-link for all
  three anomalies (E1--E3 in \texttt{eci\_apparatus\_links.md}).
  The muon $g-2$ and R(K) anomalies are in any case resolved as of 2025--2026.

\item[Neutrino masses and mass ordering.]
  No apparatus link.  ECI is silent on the neutrino sector.

\item[Dark matter direct-detection (XENONnT NR excess, 4$\sigma$).]
  The dark-dimension axiom (A5) provides a $\Delta N_{\rm eff}({\rm KK})$
  constraint, not a direct-detection cross-section.  The Bianchi VIII
  Mixmaster DM-relic route is killed by Theorem T2.

\end{description}

% ---------------------------------------------------------------------------
\subsection{Recommended Reading on What ECI Does Not Solve}
\label{subsec:reading}
% ---------------------------------------------------------------------------

Readers seeking frameworks that address the problems in
\S\ref{subsec:outside} are directed to the following.

\begin{itemize}

\item \textbf{Cosmological constant.}
  Weinberg 2000 (arXiv:astro-ph/0005265) remains the definitive review
  of why no symmetry argument succeeds.  Bousso--Polchinski 2000
  (hep-th/0004134) is the landscape mechanism.  No paper in the
  algebraic-QFT tradition (including the present one) derives
  $\rho_\Lambda / M_P^4 \approx 10^{-121}$.

\item \textbf{Hubble tension.}
  Verde--Treu--Riess 2019 (arXiv:1907.10625) review the tension.
  EDE (Poulin et al.\ 2019, arXiv:1811.04083) and NMC modifications
  (Wolf--D'Amico--Pace 2025, arXiv:2504.07679) are among the active
  explanatory approaches; Levier 1B rules out the latter at current
  sensitivity.

\item \textbf{Entropic gravity vs ECI's algebraic NO-GO.}
  Bianconi 2025 (arXiv:2502.02661) proposes an entropic origin of
  gravity.  The forthcoming ECI comment (B8 in flight) identifies the
  Lorentzian-signature gap; the comment is a mathematical objection, not
  a claim that ECI resolves dark energy.

\item \textbf{Connes--Chamseddine NCG Standard Model.}
  Chamseddine--Connes 1997 (hep-th/9606001) generates the SM action from
  the spectral action principle but does not pin $\Lambda$ or generation
  count.  The NCG approach and ECI are orthogonal: NCG addresses the SM
  action at the algebra-of-observables level; ECI addresses the
  vacuum-state structure on cosmological backgrounds.

\item \textbf{Selection-principle frameworks.}
  The Penrose--Weyl curvature hypothesis (Penrose, \emph{Cycles of Time},
  2010) and the Past Hypothesis (Albert 2000) are the standard frameworks
  for the cosmological arrow of time.  Theorem T3 of ECI gives an algebraic
  correlate, but the low-complexity boundary condition remains a posit in
  all frameworks, including ECI.

\item \textbf{CLPW observer algebras.}
  Chandrasekaran--Penington--Witten 2023 (arXiv:2206.10780) is the parent
  framework for the type-II$_\infty$ construction.  Witten 2022
  (arXiv:2112.12828) is the foundational paper on the crossed-product
  algebra in de Sitter.  These references establish what \emph{can} be
  proven at the observer-algebra level; ECI extends them to FRW backgrounds
  but does not supersede their scope.

\end{itemize}

% ---------------------------------------------------------------------------
% Local bibliography entries referenced above.
% Add these to eciNotes.bib / eci.bib if not already present.
% ---------------------------------------------------------------------------
% @misc{ECI_cosmo_const_audit,
%   author = {Remondi\`ere, Kevin},
%   title  = {Honest reality check --- cosmological constant and arrow-of-time
%             claims},
%   year   = {2026},
%   note   = {Internal audit note,
%             \texttt{notes/cosmo\_const\_aot\_honest\_2026\_05\_02.md}}
% }
 after last body section.
%
% This section documents the bounded scope of the ECI framework.
% It is NOT a defeat — it is a calibration statement.
% Every paragraph is written to pass the critic's smell test:
%   "Can a hostile referee accuse this of overclaiming?"
% If the answer is yes, the paragraph has been rewritten.

\section{Scope and Limitations of the ECI Framework}
\label{sec:scope}

The Entanglement--Complexity--Information (ECI) programme is a
mathematical-physics framework in the post-CLPW algebraic-quantum-gravity
tradition~\cite{Chandrasekaran2023,Chandrasekaran2022}.
It constructs a type-II$_\infty$ von Neumann algebra for FRW cosmological
diamonds via a quantum-reference-frame crossed product, derives a modular
(Krylov) complexity measure on that algebra, and uses the Bisognano--Wichmann
theorem and Hadamard-state theory to characterize which boundary conditions
are algebraically admissible near the Big Bang.  The framework's value is
precision and exclusion: it gives sharp structural results about
\emph{what is and is not algebraically possible} in cosmological quantum
field theory on a fixed background.

This section states, with equal emphasis, what ECI claims and what it does
not claim.  The negative results listed in \S\ref{subsec:nogo} are presented
as successes of calibration, not as failures of the programme.  A framework
that has explicitly tested and retired seventeen candidate bridges to
phenomenology has narrowed its own scope in a scientifically productive way.

% ---------------------------------------------------------------------------
\subsection{What ECI Claims and Is Positioned to Prove}
\label{subsec:claims}
% ---------------------------------------------------------------------------

\textbf{Core algebraic apparatus.}  The following results are established
within the framework, with the level of proof noted explicitly.

\begin{enumerate}

\item \textbf{FRW type-II$_\infty$ crossed product (Theorems 3.5--3.6,
  \texttt{paper/frw\_typeII\_note}).}
  The local observable algebra on a comoving past-light-cone diamond in
  radiation-era or matter-era FRW is a type-II$_\infty$ von Neumann factor
  when constructed as a quantum-reference-frame crossed product $\mathcal{A}(D_O)
  \rtimes_\sigma \mathbb{R}$.  The construction follows Witten~\cite{Witten2022},
  Chandrasekaran--Penington--Witten~\cite{Chandrasekaran2023}, and
  Leutheusser--Liu~\cite{Leutheusser2022}; the FRW adaptation is due to this
  work.  The factor carries a canonical tracial weight, but this trace fixes
  only a \emph{relative} scale (not an absolute energy density).

\item \textbf{Conformal pullback of the stress tensor (Lemma 3.3,
  sympy-verified).}
  The unitary $U$ implementing the FRW$\to$Minkowski conformal rescaling
  $a(\eta)$ satisfies $U^{-1}\,T^{\rm Mink}_{00}\,U = a^2(\eta)\,
  \mathcal{T}^{\rm FRW}_{00}$ off-shell, to all orders in the Taylor-mode
  expansion of the test-function support.  The conjugation is the
  backbone of the $K_{\rm FRW}$ closed-form derivation.

\item \textbf{$K_{\rm FRW}$ closed form (R-piste 6, $\sim$80\% paper-ready).}
  For a massless conformally-coupled scalar on a radiation-era FRW
  past-light-cone diamond $D_R$ of comoving radius $R_{\rm conf}$,
  the modular Hamiltonian is
  \begin{equation}
    K_{\rm FRW} \;=\; \frac{\pi}{R_{\rm conf}}
    \int_{D_R} d^3x\;
    \bigl[R_{\rm conf}^2 - (\eta - \eta_c)^2 - |\mathbf{x}-\mathbf{x}_c|^2\bigr]\,
    a^2(\eta)\,\mathcal{T}^{\rm FRW}_{00},
    \label{eq:kfrw}
  \end{equation}
  derived by conformal pullback from the Hislop--Longo
  Minkowski-diamond formula~\cite{HislopLongo1982}.  The prefactor
  $\pi/R_{\rm conf}$ is structural (Bisognano--Wichmann unimodular
  normalization $\beta=2\pi$, rescaled by the CHM08 Unruh relation); it is
  \emph{not} matched post-hoc to any observational value.

\item \textbf{Krylov-Diameter correspondence (Theorem 4, conditional).}
  In the late-modular-time asymptotic regime, the Krylov spread complexity
  $C_k(t)$ on the type-II$_\infty$ crossed-product algebra satisfies
  \begin{equation}
    \frac{1}{C_k(t)}\frac{dC_k}{dt} \;=\; \frac{1}{R_{\rm proper}(\eta_c)}
    + O(C_k^{-1}),
    \label{eq:krylov_rate}
  \end{equation}
  where $R_{\rm proper}(\eta_c) = a(\eta_c)\,R_{\rm conf}$ is the gauge-invariant
  proper radius of the diamond at conformal time $\eta_c$.  This identity is
  gauge-invariant (both sides are spacetime scalars) and is verified
  symbolically era-by-era (sympy).
  \textbf{Caveat:} the result is conditional on the \emph{restricted-class}
  Lanczos asymptotic $b_n = \pi n$ (Block A1, Open Question~1 in
  \texttt{krylov\_diameter.tex}) holding for the modular Hamiltonian's
  spectral class on a type-II$_\infty$ KMS state.  This transfer from the
  type-III$_1$ / CFT$_2$ setting (where the universal operator-growth
  hypothesis is established) to type-II$_\infty$ is \emph{not yet proved}.
  The conditional result is clean and publishable; the unconditional result
  requires Block A1 closure.

\item \textbf{Algebraic Weyl Curvature Hypothesis — Big-Bang
  non-Hadamard boundary (T2-FRW and T2-Bianchi I, $\sim$80\%
  paper-ready).}
  No state in the BFV--Verch Hadamard folium of $\Omega_{\rm FRW}$ extends
  to $\eta = 0$ in flat FRW (Theorem T2-FRW, three independent
  sympy-verified obstructions) or in Bianchi I (Theorem T2-Bianchi I, two
  independent IR divergences).  The physics consequence (Theorem T3) is that
  the algebras $({\mathcal A}(D_\infty), \Omega_{\rm FRW})$ and
  $(\mathcal{A}(D_{\rm BB}), \Omega_{\rm FRW})$ are not isomorphic, giving a
  rigorous algebraic correlate of the Penrose--Weyl curvature hypothesis.
  \textbf{Convention caveat:} the result is stated within the Hollands--Wald
  Hadamard microlocal convention~\cite{HollandsWald2001,BFV2003};
  alternative UV regulators may alter the IR divergence structure.
  The result does \emph{not} constitute a derivation of the
  cosmological arrow of time from first principles: the low-complexity
  initial condition (Past Hypothesis) is assumed as a boundary condition,
  not derived.

\item \textbf{Levier 1B MCMC (observational, not algebraic).}
  A joint MCMC run over Planck 2018 + ACT DR6 + DESI DR2 + Pantheon+ +
  KiDS-1000, with non-minimal coupling $\xi_\chi$ free, returns:
  $H_0 = 67.67 \pm 0.57$ km/s/Mpc, $\xi_\chi = -0.00003$
  $[-0.0164, +0.0163]$ (68\% CL), $\log B_{10} = -1.37$ (Savage--Dickey,
  weak/moderate evidence for $\Lambda$CDM over NMC).
  The ECI programme provided the NMC prior structure; the MCMC is an
  observational constraint, not an algebraic theorem.

\end{enumerate}

% ---------------------------------------------------------------------------
\subsection{What ECI Explicitly Does NOT Claim: the NO-GO Catalogue}
\label{subsec:nogo}
% ---------------------------------------------------------------------------

The following seventeen candidate connections between the ECI apparatus and
physical observables have been tested and returned negative or null results.
Each entry is a success of calibration.  They are listed here explicitly so
that future readers, referees, and automated agents do not re-open avenues
that have already been closed with stated reasons.

\subsubsection*{The five Pistes (systematic programme failures)}

\begin{description}

\item[Piste 1 (partial positive only).]
  $H = \text{Krylov rate}$: the identity
  $\tfrac{1}{C_k}\tfrac{dC_k}{dt} = H(t)$ holds exactly for a
  radiation-era Hubble-horizon diamond (gauge-invariant, sympy-verified) and
  is off by a factor of $1/2$ for a matter-era past-light-cone diamond.
  The match in the de Sitter and Hubble-horizon cases is tautological
  (the diamond is defined to have $R_{\rm proper} = H^{-1}$).
  \emph{ECI does not claim a general derivation of the Hubble parameter
  from Krylov complexity.}

\item[Piste 2 (NEGATIVE).]
  Era transitions as Jones--Longo subfactor inclusions: the algebra
  inclusion at the radiation-to-matter transition has Jones index equal to 1
  (the factors are isomorphic) and the subfactor is trivial.
  \emph{No cosmological transition amplitude is derived from subfactor theory.}

\item[Piste 3 (TAUTOLOGICAL).]
  $\Lambda$ from modular relaxation: the type-II$_\infty$ tracial weight
  fixes only a relative scale; equating it to $\rho_\Lambda / M_P^4 \approx
  10^{-121}$ requires a 121-order normalisation choice that is not derivable
  from the algebra alone.  The trace is gauge-fixed by an arbitrary modular
  normalisation, not by a physical energy density.
  \emph{ECI does not contribute a derivation of the cosmological constant.}
  (See also \S\ref{subsec:outside}.)

\item[Piste 4 (NO-GO).]
  Riemann zeros as cosmological modular Hamiltonian spectrum: the spectrum
  of $K_{\rm FRW}$ (eq.~\eqref{eq:kfrw}) is continuous, not discrete;
  the Riemann zeros form a discrete set.  A continuous-to-discrete
  mapping would require compactification not present in the apparatus.
  \emph{ECI does not offer a route to the Riemann Hypothesis via cosmology.}

\item[Piste 5 (NO-GO).]
  SM field content from modular 2-categorical compatibility: the
  type-I (finite-dimensional) algebra arising from any fixed Fock-space
  sector has only \emph{inner} modular automorphisms; the SM gauge group
  cannot be read off from the modular automorphism group of a type-I factor.
  \emph{ECI does not derive the Standard Model field content.}

\end{description}

\subsubsection*{The ten R-pistes (targeted rapid explorations)}

\begin{description}

\item[R1 (NUMEROLOGY).]
  $\Lambda$ from T1$\leftrightarrow$T2 algebraic asymmetry: the numerical
  match relied on a post-hoc normalisation; no structural reason for the
  asymmetry to equal $\rho_\Lambda$.

\item[R2 (POSITIVE BYPRODUCT, not $H_0$).]
  $H_0$ from Krylov-Diameter Theorem 4, Block A: the sympy-verified
  $\langle p_3 \rangle$ closed form is a genuine result, but it does not
  constitute a derivation of $H_0$.  The Krylov rate equals $1/R_{\rm proper}$,
  not $H_0$, in the matter era (factor of $1/2$ discrepancy;
  \S\ref{subsec:claims} item 4).

\item[R3 (KILLED: Verch 1994).]
  $S_8$ as BFV folium-projection ambiguity: Verch 1994~\cite{Verch1994}
  proves local quasi-equivalence of all Hadamard states; the alleged
  ``projection ambiguity'' does not exist within the folium.
  Levier 1B MCMC confirms null: NMC alone does not reduce the
  $S_8$ tension at Planck + KiDS-1000 sensitivity.

\item[R4 (REFUTED: Lie algebra).]
  SM three generations from Bianchi IX $SU(2)$ Kasner axes: the
  dimension count of the $\mathfrak{su}(2)$ BKL algebra is 3 (the three
  Kasner directions), but this is the rank of the spatial isometry group,
  not a fermion-generation count.  The Lie-algebraic bridge is not present.

\item[R5 (CONDITIONAL COMPANION, downgraded).]
  Riemann Hypothesis via Maass--Selberg + theta-correspondence:
  the earlier framing as an RH-proof has been retracted.
  The surviving claim is a \emph{conditional companion}: if the
  BFV-folium $\leftrightarrow$ CCM Eisenstein spectral correspondence holds,
  certain spectral conditions on the modular Hamiltonian would follow.
  No unconditional RH implication is claimed.

\item[R6 (CONDITIONAL COMPANION).]
  AdS/CFT in FRW type-II$_\infty$ frame: the original framing
  (``ECI establishes AdS/CFT for FRW'') is falsified.  The surviving
  claim is structural: the FRW type-II$_\infty$ factor is compatible
  with CLPW-style observer algebras, but the full holographic dictionary
  in an FRW background is not established.

\item[R7 (NUMEROLOGY).]
  Hierarchy from $c'$ modular trace ratios: the numerical match between
  the $c'$ parameter and $M_W/M_P$ depends on the same arbitrary
  normalisation identified in Piste 3.  No structural derivation.

\item[R8 (REFUTED: Weyl law).]
  Generation count from PSL$(2,\mathbb{Z})\backslash\mathbb{H}$ Eisenstein/cusp
  sectors: the Weyl law on the modular surface gives a
  \emph{continuous} eigenvalue density $\sim T \log T$, not a finite
  count of three.  The sector-counting route is blocked.

\item[R9 (KILLED: Theorem T2).]
  Dark matter from Bianchi VIII Mixmaster Kasner-bounce relics:
  Theorem T2-Bianchi I (and its extension to Bianchi VIII type-A)
  shows that no Hadamard state extends to the initial Kasner singularity;
  the bounce-relic mechanism requires initial conditions in the folium,
  which do not exist.

\item[R10 (NO-GO: DKN walls).]
  SM gauge from BKL $\leftrightarrow$ $E_{10}$ Kac--Moody:
  the Damour--Kleinschmidt--Nicolai billiard walls encode the
  $E_{10}$ structure but not the SM gauge factors $SU(3)\times SU(2)
  \times U(1)$; the coset decomposition does not produce SM multiplets.

\end{description}

% ---------------------------------------------------------------------------
\subsection{Open Problems Within ECI's Scope}
\label{subsec:open_within}
% ---------------------------------------------------------------------------

The following problems are within the ECI apparatus's reach and represent
the programme's live research frontier.  They are listed with their
current readiness level.

\begin{enumerate}

\item \textbf{Block A1 — unconditional UOGH transfer (open).}
  Proving $b_n = \pi n$ as an asymptotic on a type-II$_\infty$ KMS state
  without appealing to the universal operator-growth hypothesis in the
  type-III$_1$ / CFT$_2$ setting.  Completion would remove the ``conditional''
  caveat from Theorem 4.
  \emph{Status:} restricted-class Vanlessen/Levin--Lubinsky route in progress
  (B6 wave-3); full closure estimated at 2--3 months.

\item \textbf{Connes-embedding status of the FRW II$_\infty$ factor (open).}
  Is the crossed-product factor $\mathcal{A}(D_O) \rtimes_\sigma \mathbb{R}$
  approximately finite-dimensional (hyperfinite)?
  By Connes--Takesaki 1977 the answer is almost certainly yes (a crossed
  product of a hyperfinite type-III$_1$ factor by $\mathbb{R}$ is hyperfinite
  II$_\infty$), which would imply every cosmological observer in the BFV
  folium sees the same unique factor.
  \emph{Status:} 6--8 week short note; math.OA submission.

\item \textbf{T2-Bianchi V and T2-Bianchi IX extensions (open).}
  Theorem T2-Bianchi I is proved.  Extension to Bianchi V (negatively-curved
  spatial sections) is conditional on Hadamard-state existence on $H^3$;
  extension to Bianchi IX requires the BKL invariant-measure analysis of the
  Mixmaster attractor.
  \emph{Status:} partial proofs in \texttt{notes/T2\_bianchi\_2026\_05\_03\_jour2};
  3--8 weeks per case.

\item \textbf{BEC analogue-FRW experimental falsifier (open).}
  Lemma 3.3 (conformal pullback) predicts a modular spectral density on a
  BEC analogue of the FRW past-light-cone diamond.  Comparison against
  BEC entanglement-Hamiltonian tomography (Dalmonte--Vermersch--Zoller 2018
  framework) would constitute the first ECI experimental falsifier outside
  cosmological MCMC.
  \emph{Status:} 6--9 months; requires a BEC-experimental co-author.

\item \textbf{Algebraic WCH paper submission (near-ready).}
  The \texttt{paper/algebraic\_wch\_bianchi/note.tex} is 80\% paper-ready
  for Foundations of Physics.  Theorem T2-Bianchi I is the central result;
  the paper requires 2--3 weeks of polish.

\item \textbf{Bianconi 2025 PRD Comment (near-ready).}
  A 4-page note pointing out a Lorentzian-signature gap in Bianconi's
  entropic-gravity framework~\cite{Bianconi2025} has been drafted.
  The comment does \emph{not} claim ECI resolves dark energy or $\Lambda$;
  it is a structural mathematical objection
  ($\mathrm{Tr}(\tilde{g}\ln\tilde{g}^{-1}) \neq 0$ in the Lorentzian case)
  and an identification of where the ECI algebraic approach diverges from
  the entropic approach.
  \emph{Status:} 1--2 weeks to submission.

\end{enumerate}

% ---------------------------------------------------------------------------
\subsection{Open Problems Explicitly Outside ECI's Apparatus Reach}
\label{subsec:outside}
% ---------------------------------------------------------------------------

The following problems are \emph{outside} the ECI apparatus's reach.
This is not a temporary limitation to be overcome by adding more structure:
in each case a structural reason is given for why the apparatus does not
and cannot couple.

\begin{description}

\item[Cosmological constant $\Lambda \approx 10^{-122}\,M_P^4$.]
  The type-II$_\infty$ tracial weight fixes a relative scale, not an absolute
  energy density.  Equating the trace to $\rho_\Lambda / M_P^4$ introduces a
  normalisation that must be fixed from outside the algebra.
  Three independent arithmetic checks (NCG spectral cutoff, $E_8$ instanton,
  ECI internal $\kappa_R / C_k^{\rm max}$) each fail by 14--89 orders of
  magnitude~\cite{ECI_cosmo_const_audit}.
  \emph{No ECI paper claims a derivation of $\Lambda$.}

\item[Hubble tension ($H_0 = 73.04 \pm 1.04$ vs $67.67 \pm 0.57$
  km/s/Mpc, 4.5$\sigma$, UNRESOLVED).]
  Levier 1B (NMC $\xi_\chi$ free, full Planck + ACT + DESI + Pantheon+
  + KiDS-1000) returns $\log B_{10} = -1.37$ (weak/moderate evidence
  \emph{for} $\Lambda$CDM over NMC).  The NMC sector does not move $H_0$.
  Algebraically, $K_{\rm FRW}$ contains no Hubble-flow normalisation;
  Theorem 4 gives $1/R_{\rm proper}$, not $H_0$, as the Krylov rate in the
  matter era.
  \emph{ECI does not claim to resolve the Hubble tension.
  The 4.5$\sigma$ discrepancy remains open.}

\item[Standard Model particle content and gauge structure.]
  Piste 5 (inner modular automorphisms) and R4, R8, R10 (various
  generation-counting routes) all return NO-GO.  The modular automorphism
  group of any type-I sector is inner; the SM gauge factors are not
  extractable from the BFV-folium algebra.
  \emph{ECI does not predict or derive the SM gauge group, particle masses,
  or fermion generation count.}

\item[$S_8$ tension (KiDS vs Planck, $\sim$2.6$\sigma$).]
  Wolf + 2025 NMC mechanism and Levier 1B both return null on $S_8$
  (R-piste S$_8$ NEGATIVE).  Adding $\sigma_8$ as a derived parameter in
  a future MCMC run is advised before closing this entry, but NMC alone
  is not expected to resolve the tension.

\item[Swampland conjectures and string-landscape Λ-selection.]
  The Swampland $\times$ NMC cross-constraint (D8) tightens $|\xi|$ by
  $\approx 250\times$ and shows that the ECI $\chi$ field cannot also be the
  dark-dimension bulk mode (A4 $\neq$ A5).  Beyond this \emph{constraint},
  ECI does not contribute to swampland physics, de Sitter stability, or
  landscape statistics.

\item[Galaxy-scale structure (JWST early massive galaxies, $z > 10$).]
  ECI's apparatus operates on homogeneous--isotropic local diamonds.
  It has no coupling to halo abundance, star-formation efficiency, or
  first-light timing.  No ECI result bears on the JWST stellar-mass excess.

\item[SM precision observables (W-mass, muon $g-2$, lepton universality).]
  ECI's operator-algebra machinery has no coupling to loop corrections in
  perturbative QED/QCD.  The catalogue records no-apparatus-link for all
  three anomalies (E1--E3 in \texttt{eci\_apparatus\_links.md}).
  The muon $g-2$ and R(K) anomalies are in any case resolved as of 2025--2026.

\item[Neutrino masses and mass ordering.]
  No apparatus link.  ECI is silent on the neutrino sector.

\item[Dark matter direct-detection (XENONnT NR excess, 4$\sigma$).]
  The dark-dimension axiom (A5) provides a $\Delta N_{\rm eff}({\rm KK})$
  constraint, not a direct-detection cross-section.  The Bianchi VIII
  Mixmaster DM-relic route is killed by Theorem T2.

\end{description}

% ---------------------------------------------------------------------------
\subsection{Recommended Reading on What ECI Does Not Solve}
\label{subsec:reading}
% ---------------------------------------------------------------------------

Readers seeking frameworks that address the problems in
\S\ref{subsec:outside} are directed to the following.

\begin{itemize}

\item \textbf{Cosmological constant.}
  Weinberg 2000 (arXiv:astro-ph/0005265) remains the definitive review
  of why no symmetry argument succeeds.  Bousso--Polchinski 2000
  (hep-th/0004134) is the landscape mechanism.  No paper in the
  algebraic-QFT tradition (including the present one) derives
  $\rho_\Lambda / M_P^4 \approx 10^{-121}$.

\item \textbf{Hubble tension.}
  Verde--Treu--Riess 2019 (arXiv:1907.10625) review the tension.
  EDE (Poulin et al.\ 2019, arXiv:1811.04083) and NMC modifications
  (Wolf--D'Amico--Pace 2025, arXiv:2504.07679) are among the active
  explanatory approaches; Levier 1B rules out the latter at current
  sensitivity.

\item \textbf{Entropic gravity vs ECI's algebraic NO-GO.}
  Bianconi 2025 (arXiv:2502.02661) proposes an entropic origin of
  gravity.  The forthcoming ECI comment (B8 in flight) identifies the
  Lorentzian-signature gap; the comment is a mathematical objection, not
  a claim that ECI resolves dark energy.

\item \textbf{Connes--Chamseddine NCG Standard Model.}
  Chamseddine--Connes 1997 (hep-th/9606001) generates the SM action from
  the spectral action principle but does not pin $\Lambda$ or generation
  count.  The NCG approach and ECI are orthogonal: NCG addresses the SM
  action at the algebra-of-observables level; ECI addresses the
  vacuum-state structure on cosmological backgrounds.

\item \textbf{Selection-principle frameworks.}
  The Penrose--Weyl curvature hypothesis (Penrose, \emph{Cycles of Time},
  2010) and the Past Hypothesis (Albert 2000) are the standard frameworks
  for the cosmological arrow of time.  Theorem T3 of ECI gives an algebraic
  correlate, but the low-complexity boundary condition remains a posit in
  all frameworks, including ECI.

\item \textbf{CLPW observer algebras.}
  Chandrasekaran--Penington--Witten 2023 (arXiv:2206.10780) is the parent
  framework for the type-II$_\infty$ construction.  Witten 2022
  (arXiv:2112.12828) is the foundational paper on the crossed-product
  algebra in de Sitter.  These references establish what \emph{can} be
  proven at the observer-algebra level; ECI extends them to FRW backgrounds
  but does not supersede their scope.

\end{itemize}

% ---------------------------------------------------------------------------
% Local bibliography entries referenced above.
% Add these to eciNotes.bib / eci.bib if not already present.
% ---------------------------------------------------------------------------
% @misc{ECI_cosmo_const_audit,
%   author = {Remondi\`ere, Kevin},
%   title  = {Honest reality check --- cosmological constant and arrow-of-time
%             claims},
%   year   = {2026},
%   note   = {Internal audit note,
%             \texttt{notes/cosmo\_const\_aot\_honest\_2026\_05\_02.md}}
% }
 after last body section.
%
% This section documents the bounded scope of the ECI framework.
% It is NOT a defeat — it is a calibration statement.
% Every paragraph is written to pass the critic's smell test:
%   "Can a hostile referee accuse this of overclaiming?"
% If the answer is yes, the paragraph has been rewritten.

\section{Scope and Limitations of the ECI Framework}
\label{sec:scope}

The Entanglement--Complexity--Information (ECI) programme is a
mathematical-physics framework in the post-CLPW algebraic-quantum-gravity
tradition~\cite{Chandrasekaran2023,Chandrasekaran2022}.
It constructs a type-II$_\infty$ von Neumann algebra for FRW cosmological
diamonds via a quantum-reference-frame crossed product, derives a modular
(Krylov) complexity measure on that algebra, and uses the Bisognano--Wichmann
theorem and Hadamard-state theory to characterize which boundary conditions
are algebraically admissible near the Big Bang.  The framework's value is
precision and exclusion: it gives sharp structural results about
\emph{what is and is not algebraically possible} in cosmological quantum
field theory on a fixed background.

This section states, with equal emphasis, what ECI claims and what it does
not claim.  The negative results listed in \S\ref{subsec:nogo} are presented
as successes of calibration, not as failures of the programme.  A framework
that has explicitly tested and retired seventeen candidate bridges to
phenomenology has narrowed its own scope in a scientifically productive way.

% ---------------------------------------------------------------------------
\subsection{What ECI Claims and Is Positioned to Prove}
\label{subsec:claims}
% ---------------------------------------------------------------------------

\textbf{Core algebraic apparatus.}  The following results are established
within the framework, with the level of proof noted explicitly.

\begin{enumerate}

\item \textbf{FRW type-II$_\infty$ crossed product (Theorems 3.5--3.6,
  \texttt{paper/frw\_typeII\_note}).}
  The local observable algebra on a comoving past-light-cone diamond in
  radiation-era or matter-era FRW is a type-II$_\infty$ von Neumann factor
  when constructed as a quantum-reference-frame crossed product $\mathcal{A}(D_O)
  \rtimes_\sigma \mathbb{R}$.  The construction follows Witten~\cite{Witten2022},
  Chandrasekaran--Penington--Witten~\cite{Chandrasekaran2023}, and
  Leutheusser--Liu~\cite{Leutheusser2022}; the FRW adaptation is due to this
  work.  The factor carries a canonical tracial weight, but this trace fixes
  only a \emph{relative} scale (not an absolute energy density).

\item \textbf{Conformal pullback of the stress tensor (Lemma 3.3,
  sympy-verified).}
  The unitary $U$ implementing the FRW$\to$Minkowski conformal rescaling
  $a(\eta)$ satisfies $U^{-1}\,T^{\rm Mink}_{00}\,U = a^2(\eta)\,
  \mathcal{T}^{\rm FRW}_{00}$ off-shell, to all orders in the Taylor-mode
  expansion of the test-function support.  The conjugation is the
  backbone of the $K_{\rm FRW}$ closed-form derivation.

\item \textbf{$K_{\rm FRW}$ closed form (R-piste 6, $\sim$80\% paper-ready).}
  For a massless conformally-coupled scalar on a radiation-era FRW
  past-light-cone diamond $D_R$ of comoving radius $R_{\rm conf}$,
  the modular Hamiltonian is
  \begin{equation}
    K_{\rm FRW} \;=\; \frac{\pi}{R_{\rm conf}}
    \int_{D_R} d^3x\;
    \bigl[R_{\rm conf}^2 - (\eta - \eta_c)^2 - |\mathbf{x}-\mathbf{x}_c|^2\bigr]\,
    a^2(\eta)\,\mathcal{T}^{\rm FRW}_{00},
    \label{eq:kfrw}
  \end{equation}
  derived by conformal pullback from the Hislop--Longo
  Minkowski-diamond formula~\cite{HislopLongo1982}.  The prefactor
  $\pi/R_{\rm conf}$ is structural (Bisognano--Wichmann unimodular
  normalization $\beta=2\pi$, rescaled by the CHM08 Unruh relation); it is
  \emph{not} matched post-hoc to any observational value.

\item \textbf{Krylov-Diameter correspondence (Theorem 4, conditional).}
  In the late-modular-time asymptotic regime, the Krylov spread complexity
  $C_k(t)$ on the type-II$_\infty$ crossed-product algebra satisfies
  \begin{equation}
    \frac{1}{C_k(t)}\frac{dC_k}{dt} \;=\; \frac{1}{R_{\rm proper}(\eta_c)}
    + O(C_k^{-1}),
    \label{eq:krylov_rate}
  \end{equation}
  where $R_{\rm proper}(\eta_c) = a(\eta_c)\,R_{\rm conf}$ is the gauge-invariant
  proper radius of the diamond at conformal time $\eta_c$.  This identity is
  gauge-invariant (both sides are spacetime scalars) and is verified
  symbolically era-by-era (sympy).
  \textbf{Caveat:} the result is conditional on the \emph{restricted-class}
  Lanczos asymptotic $b_n = \pi n$ (Block A1, Open Question~1 in
  \texttt{krylov\_diameter.tex}) holding for the modular Hamiltonian's
  spectral class on a type-II$_\infty$ KMS state.  This transfer from the
  type-III$_1$ / CFT$_2$ setting (where the universal operator-growth
  hypothesis is established) to type-II$_\infty$ is \emph{not yet proved}.
  The conditional result is clean and publishable; the unconditional result
  requires Block A1 closure.

\item \textbf{Algebraic Weyl Curvature Hypothesis — Big-Bang
  non-Hadamard boundary (T2-FRW and T2-Bianchi I, $\sim$80\%
  paper-ready).}
  No state in the BFV--Verch Hadamard folium of $\Omega_{\rm FRW}$ extends
  to $\eta = 0$ in flat FRW (Theorem T2-FRW, three independent
  sympy-verified obstructions) or in Bianchi I (Theorem T2-Bianchi I, two
  independent IR divergences).  The physics consequence (Theorem T3) is that
  the algebras $({\mathcal A}(D_\infty), \Omega_{\rm FRW})$ and
  $(\mathcal{A}(D_{\rm BB}), \Omega_{\rm FRW})$ are not isomorphic, giving a
  rigorous algebraic correlate of the Penrose--Weyl curvature hypothesis.
  \textbf{Convention caveat:} the result is stated within the Hollands--Wald
  Hadamard microlocal convention~\cite{HollandsWald2001,BFV2003};
  alternative UV regulators may alter the IR divergence structure.
  The result does \emph{not} constitute a derivation of the
  cosmological arrow of time from first principles: the low-complexity
  initial condition (Past Hypothesis) is assumed as a boundary condition,
  not derived.

\item \textbf{Levier 1B MCMC (observational, not algebraic).}
  A joint MCMC run over Planck 2018 + ACT DR6 + DESI DR2 + Pantheon+ +
  KiDS-1000, with non-minimal coupling $\xi_\chi$ free, returns:
  $H_0 = 67.67 \pm 0.57$ km/s/Mpc, $\xi_\chi = -0.00003$
  $[-0.0164, +0.0163]$ (68\% CL), $\log B_{10} = -1.37$ (Savage--Dickey,
  weak/moderate evidence for $\Lambda$CDM over NMC).
  The ECI programme provided the NMC prior structure; the MCMC is an
  observational constraint, not an algebraic theorem.

\end{enumerate}

% ---------------------------------------------------------------------------
\subsection{What ECI Explicitly Does NOT Claim: the NO-GO Catalogue}
\label{subsec:nogo}
% ---------------------------------------------------------------------------

The following seventeen candidate connections between the ECI apparatus and
physical observables have been tested and returned negative or null results.
Each entry is a success of calibration.  They are listed here explicitly so
that future readers, referees, and automated agents do not re-open avenues
that have already been closed with stated reasons.

\subsubsection*{The five Pistes (systematic programme failures)}

\begin{description}

\item[Piste 1 (partial positive only).]
  $H = \text{Krylov rate}$: the identity
  $\tfrac{1}{C_k}\tfrac{dC_k}{dt} = H(t)$ holds exactly for a
  radiation-era Hubble-horizon diamond (gauge-invariant, sympy-verified) and
  is off by a factor of $1/2$ for a matter-era past-light-cone diamond.
  The match in the de Sitter and Hubble-horizon cases is tautological
  (the diamond is defined to have $R_{\rm proper} = H^{-1}$).
  \emph{ECI does not claim a general derivation of the Hubble parameter
  from Krylov complexity.}

\item[Piste 2 (NEGATIVE).]
  Era transitions as Jones--Longo subfactor inclusions: the algebra
  inclusion at the radiation-to-matter transition has Jones index equal to 1
  (the factors are isomorphic) and the subfactor is trivial.
  \emph{No cosmological transition amplitude is derived from subfactor theory.}

\item[Piste 3 (TAUTOLOGICAL).]
  $\Lambda$ from modular relaxation: the type-II$_\infty$ tracial weight
  fixes only a relative scale; equating it to $\rho_\Lambda / M_P^4 \approx
  10^{-121}$ requires a 121-order normalisation choice that is not derivable
  from the algebra alone.  The trace is gauge-fixed by an arbitrary modular
  normalisation, not by a physical energy density.
  \emph{ECI does not contribute a derivation of the cosmological constant.}
  (See also \S\ref{subsec:outside}.)

\item[Piste 4 (NO-GO).]
  Riemann zeros as cosmological modular Hamiltonian spectrum: the spectrum
  of $K_{\rm FRW}$ (eq.~\eqref{eq:kfrw}) is continuous, not discrete;
  the Riemann zeros form a discrete set.  A continuous-to-discrete
  mapping would require compactification not present in the apparatus.
  \emph{ECI does not offer a route to the Riemann Hypothesis via cosmology.}

\item[Piste 5 (NO-GO).]
  SM field content from modular 2-categorical compatibility: the
  type-I (finite-dimensional) algebra arising from any fixed Fock-space
  sector has only \emph{inner} modular automorphisms; the SM gauge group
  cannot be read off from the modular automorphism group of a type-I factor.
  \emph{ECI does not derive the Standard Model field content.}

\end{description}

\subsubsection*{The ten R-pistes (targeted rapid explorations)}

\begin{description}

\item[R1 (NUMEROLOGY).]
  $\Lambda$ from T1$\leftrightarrow$T2 algebraic asymmetry: the numerical
  match relied on a post-hoc normalisation; no structural reason for the
  asymmetry to equal $\rho_\Lambda$.

\item[R2 (POSITIVE BYPRODUCT, not $H_0$).]
  $H_0$ from Krylov-Diameter Theorem 4, Block A: the sympy-verified
  $\langle p_3 \rangle$ closed form is a genuine result, but it does not
  constitute a derivation of $H_0$.  The Krylov rate equals $1/R_{\rm proper}$,
  not $H_0$, in the matter era (factor of $1/2$ discrepancy;
  \S\ref{subsec:claims} item 4).

\item[R3 (KILLED: Verch 1994).]
  $S_8$ as BFV folium-projection ambiguity: Verch 1994~\cite{Verch1994}
  proves local quasi-equivalence of all Hadamard states; the alleged
  ``projection ambiguity'' does not exist within the folium.
  Levier 1B MCMC confirms null: NMC alone does not reduce the
  $S_8$ tension at Planck + KiDS-1000 sensitivity.

\item[R4 (REFUTED: Lie algebra).]
  SM three generations from Bianchi IX $SU(2)$ Kasner axes: the
  dimension count of the $\mathfrak{su}(2)$ BKL algebra is 3 (the three
  Kasner directions), but this is the rank of the spatial isometry group,
  not a fermion-generation count.  The Lie-algebraic bridge is not present.

\item[R5 (CONDITIONAL COMPANION, downgraded).]
  Riemann Hypothesis via Maass--Selberg + theta-correspondence:
  the earlier framing as an RH-proof has been retracted.
  The surviving claim is a \emph{conditional companion}: if the
  BFV-folium $\leftrightarrow$ CCM Eisenstein spectral correspondence holds,
  certain spectral conditions on the modular Hamiltonian would follow.
  No unconditional RH implication is claimed.

\item[R6 (CONDITIONAL COMPANION).]
  AdS/CFT in FRW type-II$_\infty$ frame: the original framing
  (``ECI establishes AdS/CFT for FRW'') is falsified.  The surviving
  claim is structural: the FRW type-II$_\infty$ factor is compatible
  with CLPW-style observer algebras, but the full holographic dictionary
  in an FRW background is not established.

\item[R7 (NUMEROLOGY).]
  Hierarchy from $c'$ modular trace ratios: the numerical match between
  the $c'$ parameter and $M_W/M_P$ depends on the same arbitrary
  normalisation identified in Piste 3.  No structural derivation.

\item[R8 (REFUTED: Weyl law).]
  Generation count from PSL$(2,\mathbb{Z})\backslash\mathbb{H}$ Eisenstein/cusp
  sectors: the Weyl law on the modular surface gives a
  \emph{continuous} eigenvalue density $\sim T \log T$, not a finite
  count of three.  The sector-counting route is blocked.

\item[R9 (KILLED: Theorem T2).]
  Dark matter from Bianchi VIII Mixmaster Kasner-bounce relics:
  Theorem T2-Bianchi I (and its extension to Bianchi VIII type-A)
  shows that no Hadamard state extends to the initial Kasner singularity;
  the bounce-relic mechanism requires initial conditions in the folium,
  which do not exist.

\item[R10 (NO-GO: DKN walls).]
  SM gauge from BKL $\leftrightarrow$ $E_{10}$ Kac--Moody:
  the Damour--Kleinschmidt--Nicolai billiard walls encode the
  $E_{10}$ structure but not the SM gauge factors $SU(3)\times SU(2)
  \times U(1)$; the coset decomposition does not produce SM multiplets.

\end{description}

% ---------------------------------------------------------------------------
\subsection{Open Problems Within ECI's Scope}
\label{subsec:open_within}
% ---------------------------------------------------------------------------

The following problems are within the ECI apparatus's reach and represent
the programme's live research frontier.  They are listed with their
current readiness level.

\begin{enumerate}

\item \textbf{Block A1 — unconditional UOGH transfer (open).}
  Proving $b_n = \pi n$ as an asymptotic on a type-II$_\infty$ KMS state
  without appealing to the universal operator-growth hypothesis in the
  type-III$_1$ / CFT$_2$ setting.  Completion would remove the ``conditional''
  caveat from Theorem 4.
  \emph{Status:} restricted-class Vanlessen/Levin--Lubinsky route in progress
  (B6 wave-3); full closure estimated at 2--3 months.

\item \textbf{Connes-embedding status of the FRW II$_\infty$ factor (open).}
  Is the crossed-product factor $\mathcal{A}(D_O) \rtimes_\sigma \mathbb{R}$
  approximately finite-dimensional (hyperfinite)?
  By Connes--Takesaki 1977 the answer is almost certainly yes (a crossed
  product of a hyperfinite type-III$_1$ factor by $\mathbb{R}$ is hyperfinite
  II$_\infty$), which would imply every cosmological observer in the BFV
  folium sees the same unique factor.
  \emph{Status:} 6--8 week short note; math.OA submission.

\item \textbf{T2-Bianchi V and T2-Bianchi IX extensions (open).}
  Theorem T2-Bianchi I is proved.  Extension to Bianchi V (negatively-curved
  spatial sections) is conditional on Hadamard-state existence on $H^3$;
  extension to Bianchi IX requires the BKL invariant-measure analysis of the
  Mixmaster attractor.
  \emph{Status:} partial proofs in \texttt{notes/T2\_bianchi\_2026\_05\_03\_jour2};
  3--8 weeks per case.

\item \textbf{BEC analogue-FRW experimental falsifier (open).}
  Lemma 3.3 (conformal pullback) predicts a modular spectral density on a
  BEC analogue of the FRW past-light-cone diamond.  Comparison against
  BEC entanglement-Hamiltonian tomography (Dalmonte--Vermersch--Zoller 2018
  framework) would constitute the first ECI experimental falsifier outside
  cosmological MCMC.
  \emph{Status:} 6--9 months; requires a BEC-experimental co-author.

\item \textbf{Algebraic WCH paper submission (near-ready).}
  The \texttt{paper/algebraic\_wch\_bianchi/note.tex} is 80\% paper-ready
  for Foundations of Physics.  Theorem T2-Bianchi I is the central result;
  the paper requires 2--3 weeks of polish.

\item \textbf{Bianconi 2025 PRD Comment (near-ready).}
  A 4-page note pointing out a Lorentzian-signature gap in Bianconi's
  entropic-gravity framework~\cite{Bianconi2025} has been drafted.
  The comment does \emph{not} claim ECI resolves dark energy or $\Lambda$;
  it is a structural mathematical objection
  ($\mathrm{Tr}(\tilde{g}\ln\tilde{g}^{-1}) \neq 0$ in the Lorentzian case)
  and an identification of where the ECI algebraic approach diverges from
  the entropic approach.
  \emph{Status:} 1--2 weeks to submission.

\end{enumerate}

% ---------------------------------------------------------------------------
\subsection{Open Problems Explicitly Outside ECI's Apparatus Reach}
\label{subsec:outside}
% ---------------------------------------------------------------------------

The following problems are \emph{outside} the ECI apparatus's reach.
This is not a temporary limitation to be overcome by adding more structure:
in each case a structural reason is given for why the apparatus does not
and cannot couple.

\begin{description}

\item[Cosmological constant $\Lambda \approx 10^{-122}\,M_P^4$.]
  The type-II$_\infty$ tracial weight fixes a relative scale, not an absolute
  energy density.  Equating the trace to $\rho_\Lambda / M_P^4$ introduces a
  normalisation that must be fixed from outside the algebra.
  Three independent arithmetic checks (NCG spectral cutoff, $E_8$ instanton,
  ECI internal $\kappa_R / C_k^{\rm max}$) each fail by 14--89 orders of
  magnitude~\cite{ECI_cosmo_const_audit}.
  \emph{No ECI paper claims a derivation of $\Lambda$.}

\item[Hubble tension ($H_0 = 73.04 \pm 1.04$ vs $67.67 \pm 0.57$
  km/s/Mpc, 4.5$\sigma$, UNRESOLVED).]
  Levier 1B (NMC $\xi_\chi$ free, full Planck + ACT + DESI + Pantheon+
  + KiDS-1000) returns $\log B_{10} = -1.37$ (weak/moderate evidence
  \emph{for} $\Lambda$CDM over NMC).  The NMC sector does not move $H_0$.
  Algebraically, $K_{\rm FRW}$ contains no Hubble-flow normalisation;
  Theorem 4 gives $1/R_{\rm proper}$, not $H_0$, as the Krylov rate in the
  matter era.
  \emph{ECI does not claim to resolve the Hubble tension.
  The 4.5$\sigma$ discrepancy remains open.}

\item[Standard Model particle content and gauge structure.]
  Piste 5 (inner modular automorphisms) and R4, R8, R10 (various
  generation-counting routes) all return NO-GO.  The modular automorphism
  group of any type-I sector is inner; the SM gauge factors are not
  extractable from the BFV-folium algebra.
  \emph{ECI does not predict or derive the SM gauge group, particle masses,
  or fermion generation count.}

\item[$S_8$ tension (KiDS vs Planck, $\sim$2.6$\sigma$).]
  Wolf + 2025 NMC mechanism and Levier 1B both return null on $S_8$
  (R-piste S$_8$ NEGATIVE).  Adding $\sigma_8$ as a derived parameter in
  a future MCMC run is advised before closing this entry, but NMC alone
  is not expected to resolve the tension.

\item[Swampland conjectures and string-landscape Λ-selection.]
  The Swampland $\times$ NMC cross-constraint (D8) tightens $|\xi|$ by
  $\approx 250\times$ and shows that the ECI $\chi$ field cannot also be the
  dark-dimension bulk mode (A4 $\neq$ A5).  Beyond this \emph{constraint},
  ECI does not contribute to swampland physics, de Sitter stability, or
  landscape statistics.

\item[Galaxy-scale structure (JWST early massive galaxies, $z > 10$).]
  ECI's apparatus operates on homogeneous--isotropic local diamonds.
  It has no coupling to halo abundance, star-formation efficiency, or
  first-light timing.  No ECI result bears on the JWST stellar-mass excess.

\item[SM precision observables (W-mass, muon $g-2$, lepton universality).]
  ECI's operator-algebra machinery has no coupling to loop corrections in
  perturbative QED/QCD.  The catalogue records no-apparatus-link for all
  three anomalies (E1--E3 in \texttt{eci\_apparatus\_links.md}).
  The muon $g-2$ and R(K) anomalies are in any case resolved as of 2025--2026.

\item[Neutrino masses and mass ordering.]
  No apparatus link.  ECI is silent on the neutrino sector.

\item[Dark matter direct-detection (XENONnT NR excess, 4$\sigma$).]
  The dark-dimension axiom (A5) provides a $\Delta N_{\rm eff}({\rm KK})$
  constraint, not a direct-detection cross-section.  The Bianchi VIII
  Mixmaster DM-relic route is killed by Theorem T2.

\end{description}

% ---------------------------------------------------------------------------
\subsection{Recommended Reading on What ECI Does Not Solve}
\label{subsec:reading}
% ---------------------------------------------------------------------------

Readers seeking frameworks that address the problems in
\S\ref{subsec:outside} are directed to the following.

\begin{itemize}

\item \textbf{Cosmological constant.}
  Weinberg 2000 (arXiv:astro-ph/0005265) remains the definitive review
  of why no symmetry argument succeeds.  Bousso--Polchinski 2000
  (hep-th/0004134) is the landscape mechanism.  No paper in the
  algebraic-QFT tradition (including the present one) derives
  $\rho_\Lambda / M_P^4 \approx 10^{-121}$.

\item \textbf{Hubble tension.}
  Verde--Treu--Riess 2019 (arXiv:1907.10625) review the tension.
  EDE (Poulin et al.\ 2019, arXiv:1811.04083) and NMC modifications
  (Wolf--D'Amico--Pace 2025, arXiv:2504.07679) are among the active
  explanatory approaches; Levier 1B rules out the latter at current
  sensitivity.

\item \textbf{Entropic gravity vs ECI's algebraic NO-GO.}
  Bianconi 2025 (arXiv:2502.02661) proposes an entropic origin of
  gravity.  The forthcoming ECI comment (B8 in flight) identifies the
  Lorentzian-signature gap; the comment is a mathematical objection, not
  a claim that ECI resolves dark energy.

\item \textbf{Connes--Chamseddine NCG Standard Model.}
  Chamseddine--Connes 1997 (hep-th/9606001) generates the SM action from
  the spectral action principle but does not pin $\Lambda$ or generation
  count.  The NCG approach and ECI are orthogonal: NCG addresses the SM
  action at the algebra-of-observables level; ECI addresses the
  vacuum-state structure on cosmological backgrounds.

\item \textbf{Selection-principle frameworks.}
  The Penrose--Weyl curvature hypothesis (Penrose, \emph{Cycles of Time},
  2010) and the Past Hypothesis (Albert 2000) are the standard frameworks
  for the cosmological arrow of time.  Theorem T3 of ECI gives an algebraic
  correlate, but the low-complexity boundary condition remains a posit in
  all frameworks, including ECI.

\item \textbf{CLPW observer algebras.}
  Chandrasekaran--Penington--Witten 2023 (arXiv:2206.10780) is the parent
  framework for the type-II$_\infty$ construction.  Witten 2022
  (arXiv:2112.12828) is the foundational paper on the crossed-product
  algebra in de Sitter.  These references establish what \emph{can} be
  proven at the observer-algebra level; ECI extends them to FRW backgrounds
  but does not supersede their scope.

\end{itemize}

% ---------------------------------------------------------------------------
% Local bibliography entries referenced above.
% Add these to eciNotes.bib / eci.bib if not already present.
% ---------------------------------------------------------------------------
% @misc{ECI_cosmo_const_audit,
%   author = {Remondi\`ere, Kevin},
%   title  = {Honest reality check --- cosmological constant and arrow-of-time
%             claims},
%   year   = {2026},
%   note   = {Internal audit note,
%             \texttt{notes/cosmo\_const\_aot\_honest\_2026\_05\_02.md}}
% }
 after last body section.
%
% This section documents the bounded scope of the ECI framework.
% It is NOT a defeat — it is a calibration statement.
% Every paragraph is written to pass the critic's smell test:
%   "Can a hostile referee accuse this of overclaiming?"
% If the answer is yes, the paragraph has been rewritten.

\section{Scope and Limitations of the ECI Framework}
\label{sec:scope}

The Entanglement--Complexity--Information (ECI) programme is a
mathematical-physics framework in the post-CLPW algebraic-quantum-gravity
tradition~\cite{Chandrasekaran2023,Chandrasekaran2022}.
It constructs a type-II$_\infty$ von Neumann algebra for FRW cosmological
diamonds via a quantum-reference-frame crossed product, derives a modular
(Krylov) complexity measure on that algebra, and uses the Bisognano--Wichmann
theorem and Hadamard-state theory to characterize which boundary conditions
are algebraically admissible near the Big Bang.  The framework's value is
precision and exclusion: it gives sharp structural results about
\emph{what is and is not algebraically possible} in cosmological quantum
field theory on a fixed background.

This section states, with equal emphasis, what ECI claims and what it does
not claim.  The negative results listed in \S\ref{subsec:nogo} are presented
as successes of calibration, not as failures of the programme.  A framework
that has explicitly tested and retired seventeen candidate bridges to
phenomenology has narrowed its own scope in a scientifically productive way.

% ---------------------------------------------------------------------------
\subsection{What ECI Claims and Is Positioned to Prove}
\label{subsec:claims}
% ---------------------------------------------------------------------------

\textbf{Core algebraic apparatus.}  The following results are established
within the framework, with the level of proof noted explicitly.

\begin{enumerate}

\item \textbf{FRW type-II$_\infty$ crossed product (Theorems 3.5--3.6,
  \texttt{paper/frw\_typeII\_note}).}
  The local observable algebra on a comoving past-light-cone diamond in
  radiation-era or matter-era FRW is a type-II$_\infty$ von Neumann factor
  when constructed as a quantum-reference-frame crossed product $\mathcal{A}(D_O)
  \rtimes_\sigma \mathbb{R}$.  The construction follows Witten~\cite{Witten2022},
  Chandrasekaran--Penington--Witten~\cite{Chandrasekaran2023}, and
  Leutheusser--Liu~\cite{Leutheusser2022}; the FRW adaptation is due to this
  work.  The factor carries a canonical tracial weight, but this trace fixes
  only a \emph{relative} scale (not an absolute energy density).

\item \textbf{Conformal pullback of the stress tensor (Lemma 3.3,
  sympy-verified).}
  The unitary $U$ implementing the FRW$\to$Minkowski conformal rescaling
  $a(\eta)$ satisfies $U^{-1}\,T^{\rm Mink}_{00}\,U = a^2(\eta)\,
  \mathcal{T}^{\rm FRW}_{00}$ off-shell, to all orders in the Taylor-mode
  expansion of the test-function support.  The conjugation is the
  backbone of the $K_{\rm FRW}$ closed-form derivation.

\item \textbf{$K_{\rm FRW}$ closed form (R-piste 6, $\sim$80\% paper-ready).}
  For a massless conformally-coupled scalar on a radiation-era FRW
  past-light-cone diamond $D_R$ of comoving radius $R_{\rm conf}$,
  the modular Hamiltonian is
  \begin{equation}
    K_{\rm FRW} \;=\; \frac{\pi}{R_{\rm conf}}
    \int_{D_R} d^3x\;
    \bigl[R_{\rm conf}^2 - (\eta - \eta_c)^2 - |\mathbf{x}-\mathbf{x}_c|^2\bigr]\,
    a^2(\eta)\,\mathcal{T}^{\rm FRW}_{00},
    \label{eq:kfrw}
  \end{equation}
  derived by conformal pullback from the Hislop--Longo
  Minkowski-diamond formula~\cite{HislopLongo1982}.  The prefactor
  $\pi/R_{\rm conf}$ is structural (Bisognano--Wichmann unimodular
  normalization $\beta=2\pi$, rescaled by the CHM08 Unruh relation); it is
  \emph{not} matched post-hoc to any observational value.

\item \textbf{Krylov-Diameter correspondence (Theorem 4, conditional).}
  In the late-modular-time asymptotic regime, the Krylov spread complexity
  $C_k(t)$ on the type-II$_\infty$ crossed-product algebra satisfies
  \begin{equation}
    \frac{1}{C_k(t)}\frac{dC_k}{dt} \;=\; \frac{1}{R_{\rm proper}(\eta_c)}
    + O(C_k^{-1}),
    \label{eq:krylov_rate}
  \end{equation}
  where $R_{\rm proper}(\eta_c) = a(\eta_c)\,R_{\rm conf}$ is the gauge-invariant
  proper radius of the diamond at conformal time $\eta_c$.  This identity is
  gauge-invariant (both sides are spacetime scalars) and is verified
  symbolically era-by-era (sympy).
  \textbf{Caveat:} the result is conditional on the \emph{restricted-class}
  Lanczos asymptotic $b_n = \pi n$ (Block A1, Open Question~1 in
  \texttt{krylov\_diameter.tex}) holding for the modular Hamiltonian's
  spectral class on a type-II$_\infty$ KMS state.  This transfer from the
  type-III$_1$ / CFT$_2$ setting (where the universal operator-growth
  hypothesis is established) to type-II$_\infty$ is \emph{not yet proved}.
  The conditional result is clean and publishable; the unconditional result
  requires Block A1 closure.

\item \textbf{Algebraic Weyl Curvature Hypothesis — Big-Bang
  non-Hadamard boundary (T2-FRW and T2-Bianchi I, $\sim$80\%
  paper-ready).}
  No state in the BFV--Verch Hadamard folium of $\Omega_{\rm FRW}$ extends
  to $\eta = 0$ in flat FRW (Theorem T2-FRW, three independent
  sympy-verified obstructions) or in Bianchi I (Theorem T2-Bianchi I, two
  independent IR divergences).  The physics consequence (Theorem T3) is that
  the algebras $({\mathcal A}(D_\infty), \Omega_{\rm FRW})$ and
  $(\mathcal{A}(D_{\rm BB}), \Omega_{\rm FRW})$ are not isomorphic, giving a
  rigorous algebraic correlate of the Penrose--Weyl curvature hypothesis.
  \textbf{Convention caveat:} the result is stated within the Hollands--Wald
  Hadamard microlocal convention~\cite{HollandsWald2001,BFV2003};
  alternative UV regulators may alter the IR divergence structure.
  The result does \emph{not} constitute a derivation of the
  cosmological arrow of time from first principles: the low-complexity
  initial condition (Past Hypothesis) is assumed as a boundary condition,
  not derived.

\item \textbf{Levier 1B MCMC (observational, not algebraic).}
  A joint MCMC run over Planck 2018 + ACT DR6 + DESI DR2 + Pantheon+ +
  KiDS-1000, with non-minimal coupling $\xi_\chi$ free, returns:
  $H_0 = 67.67 \pm 0.57$ km/s/Mpc, $\xi_\chi = -0.00003$
  $[-0.0164, +0.0163]$ (68\% CL), $\log B_{10} = -1.37$ (Savage--Dickey,
  weak/moderate evidence for $\Lambda$CDM over NMC).
  The ECI programme provided the NMC prior structure; the MCMC is an
  observational constraint, not an algebraic theorem.

\end{enumerate}

% ---------------------------------------------------------------------------
\subsection{What ECI Explicitly Does NOT Claim: the NO-GO Catalogue}
\label{subsec:nogo}
% ---------------------------------------------------------------------------

The following seventeen candidate connections between the ECI apparatus and
physical observables have been tested and returned negative or null results.
Each entry is a success of calibration.  They are listed here explicitly so
that future readers, referees, and automated agents do not re-open avenues
that have already been closed with stated reasons.

\subsubsection*{The five Pistes (systematic programme failures)}

\begin{description}

\item[Piste 1 (partial positive only).]
  $H = \text{Krylov rate}$: the identity
  $\tfrac{1}{C_k}\tfrac{dC_k}{dt} = H(t)$ holds exactly for a
  radiation-era Hubble-horizon diamond (gauge-invariant, sympy-verified) and
  is off by a factor of $1/2$ for a matter-era past-light-cone diamond.
  The match in the de Sitter and Hubble-horizon cases is tautological
  (the diamond is defined to have $R_{\rm proper} = H^{-1}$).
  \emph{ECI does not claim a general derivation of the Hubble parameter
  from Krylov complexity.}

\item[Piste 2 (NEGATIVE).]
  Era transitions as Jones--Longo subfactor inclusions: the algebra
  inclusion at the radiation-to-matter transition has Jones index equal to 1
  (the factors are isomorphic) and the subfactor is trivial.
  \emph{No cosmological transition amplitude is derived from subfactor theory.}

\item[Piste 3 (TAUTOLOGICAL).]
  $\Lambda$ from modular relaxation: the type-II$_\infty$ tracial weight
  fixes only a relative scale; equating it to $\rho_\Lambda / M_P^4 \approx
  10^{-121}$ requires a 121-order normalisation choice that is not derivable
  from the algebra alone.  The trace is gauge-fixed by an arbitrary modular
  normalisation, not by a physical energy density.
  \emph{ECI does not contribute a derivation of the cosmological constant.}
  (See also \S\ref{subsec:outside}.)

\item[Piste 4 (NO-GO).]
  Riemann zeros as cosmological modular Hamiltonian spectrum: the spectrum
  of $K_{\rm FRW}$ (eq.~\eqref{eq:kfrw}) is continuous, not discrete;
  the Riemann zeros form a discrete set.  A continuous-to-discrete
  mapping would require compactification not present in the apparatus.
  \emph{ECI does not offer a route to the Riemann Hypothesis via cosmology.}

\item[Piste 5 (NO-GO).]
  SM field content from modular 2-categorical compatibility: the
  type-I (finite-dimensional) algebra arising from any fixed Fock-space
  sector has only \emph{inner} modular automorphisms; the SM gauge group
  cannot be read off from the modular automorphism group of a type-I factor.
  \emph{ECI does not derive the Standard Model field content.}

\end{description}

\subsubsection*{The ten R-pistes (targeted rapid explorations)}

\begin{description}

\item[R1 (NUMEROLOGY).]
  $\Lambda$ from T1$\leftrightarrow$T2 algebraic asymmetry: the numerical
  match relied on a post-hoc normalisation; no structural reason for the
  asymmetry to equal $\rho_\Lambda$.

\item[R2 (POSITIVE BYPRODUCT, not $H_0$).]
  $H_0$ from Krylov-Diameter Theorem 4, Block A: the sympy-verified
  $\langle p_3 \rangle$ closed form is a genuine result, but it does not
  constitute a derivation of $H_0$.  The Krylov rate equals $1/R_{\rm proper}$,
  not $H_0$, in the matter era (factor of $1/2$ discrepancy;
  \S\ref{subsec:claims} item 4).

\item[R3 (KILLED: Verch 1994).]
  $S_8$ as BFV folium-projection ambiguity: Verch 1994~\cite{Verch1994}
  proves local quasi-equivalence of all Hadamard states; the alleged
  ``projection ambiguity'' does not exist within the folium.
  Levier 1B MCMC confirms null: NMC alone does not reduce the
  $S_8$ tension at Planck + KiDS-1000 sensitivity.

\item[R4 (REFUTED: Lie algebra).]
  SM three generations from Bianchi IX $SU(2)$ Kasner axes: the
  dimension count of the $\mathfrak{su}(2)$ BKL algebra is 3 (the three
  Kasner directions), but this is the rank of the spatial isometry group,
  not a fermion-generation count.  The Lie-algebraic bridge is not present.

\item[R5 (CONDITIONAL COMPANION, downgraded).]
  Riemann Hypothesis via Maass--Selberg + theta-correspondence:
  the earlier framing as an RH-proof has been retracted.
  The surviving claim is a \emph{conditional companion}: if the
  BFV-folium $\leftrightarrow$ CCM Eisenstein spectral correspondence holds,
  certain spectral conditions on the modular Hamiltonian would follow.
  No unconditional RH implication is claimed.

\item[R6 (CONDITIONAL COMPANION).]
  AdS/CFT in FRW type-II$_\infty$ frame: the original framing
  (``ECI establishes AdS/CFT for FRW'') is falsified.  The surviving
  claim is structural: the FRW type-II$_\infty$ factor is compatible
  with CLPW-style observer algebras, but the full holographic dictionary
  in an FRW background is not established.

\item[R7 (NUMEROLOGY).]
  Hierarchy from $c'$ modular trace ratios: the numerical match between
  the $c'$ parameter and $M_W/M_P$ depends on the same arbitrary
  normalisation identified in Piste 3.  No structural derivation.

\item[R8 (REFUTED: Weyl law).]
  Generation count from PSL$(2,\mathbb{Z})\backslash\mathbb{H}$ Eisenstein/cusp
  sectors: the Weyl law on the modular surface gives a
  \emph{continuous} eigenvalue density $\sim T \log T$, not a finite
  count of three.  The sector-counting route is blocked.

\item[R9 (KILLED: Theorem T2).]
  Dark matter from Bianchi VIII Mixmaster Kasner-bounce relics:
  Theorem T2-Bianchi I (and its extension to Bianchi VIII type-A)
  shows that no Hadamard state extends to the initial Kasner singularity;
  the bounce-relic mechanism requires initial conditions in the folium,
  which do not exist.

\item[R10 (NO-GO: DKN walls).]
  SM gauge from BKL $\leftrightarrow$ $E_{10}$ Kac--Moody:
  the Damour--Kleinschmidt--Nicolai billiard walls encode the
  $E_{10}$ structure but not the SM gauge factors $SU(3)\times SU(2)
  \times U(1)$; the coset decomposition does not produce SM multiplets.

\end{description}

% ---------------------------------------------------------------------------
\subsection{Open Problems Within ECI's Scope}
\label{subsec:open_within}
% ---------------------------------------------------------------------------

The following problems are within the ECI apparatus's reach and represent
the programme's live research frontier.  They are listed with their
current readiness level.

\begin{enumerate}

\item \textbf{Block A1 — unconditional UOGH transfer (open).}
  Proving $b_n = \pi n$ as an asymptotic on a type-II$_\infty$ KMS state
  without appealing to the universal operator-growth hypothesis in the
  type-III$_1$ / CFT$_2$ setting.  Completion would remove the ``conditional''
  caveat from Theorem 4.
  \emph{Status:} restricted-class Vanlessen/Levin--Lubinsky route in progress
  (B6 wave-3); full closure estimated at 2--3 months.

\item \textbf{Connes-embedding status of the FRW II$_\infty$ factor (open).}
  Is the crossed-product factor $\mathcal{A}(D_O) \rtimes_\sigma \mathbb{R}$
  approximately finite-dimensional (hyperfinite)?
  By Connes--Takesaki 1977 the answer is almost certainly yes (a crossed
  product of a hyperfinite type-III$_1$ factor by $\mathbb{R}$ is hyperfinite
  II$_\infty$), which would imply every cosmological observer in the BFV
  folium sees the same unique factor.
  \emph{Status:} 6--8 week short note; math.OA submission.

\item \textbf{T2-Bianchi V and T2-Bianchi IX extensions (open).}
  Theorem T2-Bianchi I is proved.  Extension to Bianchi V (negatively-curved
  spatial sections) is conditional on Hadamard-state existence on $H^3$;
  extension to Bianchi IX requires the BKL invariant-measure analysis of the
  Mixmaster attractor.
  \emph{Status:} partial proofs in \texttt{notes/T2\_bianchi\_2026\_05\_03\_jour2};
  3--8 weeks per case.

\item \textbf{BEC analogue-FRW experimental falsifier (open).}
  Lemma 3.3 (conformal pullback) predicts a modular spectral density on a
  BEC analogue of the FRW past-light-cone diamond.  Comparison against
  BEC entanglement-Hamiltonian tomography (Dalmonte--Vermersch--Zoller 2018
  framework) would constitute the first ECI experimental falsifier outside
  cosmological MCMC.
  \emph{Status:} 6--9 months; requires a BEC-experimental co-author.

\item \textbf{Algebraic WCH paper submission (near-ready).}
  The \texttt{paper/algebraic\_wch\_bianchi/note.tex} is 80\% paper-ready
  for Foundations of Physics.  Theorem T2-Bianchi I is the central result;
  the paper requires 2--3 weeks of polish.

\item \textbf{Bianconi 2025 PRD Comment (near-ready).}
  A 4-page note pointing out a Lorentzian-signature gap in Bianconi's
  entropic-gravity framework~\cite{Bianconi2025} has been drafted.
  The comment does \emph{not} claim ECI resolves dark energy or $\Lambda$;
  it is a structural mathematical objection
  ($\mathrm{Tr}(\tilde{g}\ln\tilde{g}^{-1}) \neq 0$ in the Lorentzian case)
  and an identification of where the ECI algebraic approach diverges from
  the entropic approach.
  \emph{Status:} 1--2 weeks to submission.

\end{enumerate}

% ---------------------------------------------------------------------------
\subsection{Open Problems Explicitly Outside ECI's Apparatus Reach}
\label{subsec:outside}
% ---------------------------------------------------------------------------

The following problems are \emph{outside} the ECI apparatus's reach.
This is not a temporary limitation to be overcome by adding more structure:
in each case a structural reason is given for why the apparatus does not
and cannot couple.

\begin{description}

\item[Cosmological constant $\Lambda \approx 10^{-122}\,M_P^4$.]
  The type-II$_\infty$ tracial weight fixes a relative scale, not an absolute
  energy density.  Equating the trace to $\rho_\Lambda / M_P^4$ introduces a
  normalisation that must be fixed from outside the algebra.
  Three independent arithmetic checks (NCG spectral cutoff, $E_8$ instanton,
  ECI internal $\kappa_R / C_k^{\rm max}$) each fail by 14--89 orders of
  magnitude~\cite{ECI_cosmo_const_audit}.
  \emph{No ECI paper claims a derivation of $\Lambda$.}

\item[Hubble tension ($H_0 = 73.04 \pm 1.04$ vs $67.67 \pm 0.57$
  km/s/Mpc, 4.5$\sigma$, UNRESOLVED).]
  Levier 1B (NMC $\xi_\chi$ free, full Planck + ACT + DESI + Pantheon+
  + KiDS-1000) returns $\log B_{10} = -1.37$ (weak/moderate evidence
  \emph{for} $\Lambda$CDM over NMC).  The NMC sector does not move $H_0$.
  Algebraically, $K_{\rm FRW}$ contains no Hubble-flow normalisation;
  Theorem 4 gives $1/R_{\rm proper}$, not $H_0$, as the Krylov rate in the
  matter era.
  \emph{ECI does not claim to resolve the Hubble tension.
  The 4.5$\sigma$ discrepancy remains open.}

\item[Standard Model particle content and gauge structure.]
  Piste 5 (inner modular automorphisms) and R4, R8, R10 (various
  generation-counting routes) all return NO-GO.  The modular automorphism
  group of any type-I sector is inner; the SM gauge factors are not
  extractable from the BFV-folium algebra.
  \emph{ECI does not predict or derive the SM gauge group, particle masses,
  or fermion generation count.}

\item[$S_8$ tension (KiDS vs Planck, $\sim$2.6$\sigma$).]
  Wolf + 2025 NMC mechanism and Levier 1B both return null on $S_8$
  (R-piste S$_8$ NEGATIVE).  Adding $\sigma_8$ as a derived parameter in
  a future MCMC run is advised before closing this entry, but NMC alone
  is not expected to resolve the tension.

\item[Swampland conjectures and string-landscape Λ-selection.]
  The Swampland $\times$ NMC cross-constraint (D8) tightens $|\xi|$ by
  $\approx 250\times$ and shows that the ECI $\chi$ field cannot also be the
  dark-dimension bulk mode (A4 $\neq$ A5).  Beyond this \emph{constraint},
  ECI does not contribute to swampland physics, de Sitter stability, or
  landscape statistics.

\item[Galaxy-scale structure (JWST early massive galaxies, $z > 10$).]
  ECI's apparatus operates on homogeneous--isotropic local diamonds.
  It has no coupling to halo abundance, star-formation efficiency, or
  first-light timing.  No ECI result bears on the JWST stellar-mass excess.

\item[SM precision observables (W-mass, muon $g-2$, lepton universality).]
  ECI's operator-algebra machinery has no coupling to loop corrections in
  perturbative QED/QCD.  The catalogue records no-apparatus-link for all
  three anomalies (E1--E3 in \texttt{eci\_apparatus\_links.md}).
  The muon $g-2$ and R(K) anomalies are in any case resolved as of 2025--2026.

\item[Neutrino masses and mass ordering.]
  No apparatus link.  ECI is silent on the neutrino sector.

\item[Dark matter direct-detection (XENONnT NR excess, 4$\sigma$).]
  The dark-dimension axiom (A5) provides a $\Delta N_{\rm eff}({\rm KK})$
  constraint, not a direct-detection cross-section.  The Bianchi VIII
  Mixmaster DM-relic route is killed by Theorem T2.

\end{description}

% ---------------------------------------------------------------------------
\subsection{Recommended Reading on What ECI Does Not Solve}
\label{subsec:reading}
% ---------------------------------------------------------------------------

Readers seeking frameworks that address the problems in
\S\ref{subsec:outside} are directed to the following.

\begin{itemize}

\item \textbf{Cosmological constant.}
  Weinberg 2000 (arXiv:astro-ph/0005265) remains the definitive review
  of why no symmetry argument succeeds.  Bousso--Polchinski 2000
  (hep-th/0004134) is the landscape mechanism.  No paper in the
  algebraic-QFT tradition (including the present one) derives
  $\rho_\Lambda / M_P^4 \approx 10^{-121}$.

\item \textbf{Hubble tension.}
  Verde--Treu--Riess 2019 (arXiv:1907.10625) review the tension.
  EDE (Poulin et al.\ 2019, arXiv:1811.04083) and NMC modifications
  (Wolf--D'Amico--Pace 2025, arXiv:2504.07679) are among the active
  explanatory approaches; Levier 1B rules out the latter at current
  sensitivity.

\item \textbf{Entropic gravity vs ECI's algebraic NO-GO.}
  Bianconi 2025 (arXiv:2502.02661) proposes an entropic origin of
  gravity.  The forthcoming ECI comment (B8 in flight) identifies the
  Lorentzian-signature gap; the comment is a mathematical objection, not
  a claim that ECI resolves dark energy.

\item \textbf{Connes--Chamseddine NCG Standard Model.}
  Chamseddine--Connes 1997 (hep-th/9606001) generates the SM action from
  the spectral action principle but does not pin $\Lambda$ or generation
  count.  The NCG approach and ECI are orthogonal: NCG addresses the SM
  action at the algebra-of-observables level; ECI addresses the
  vacuum-state structure on cosmological backgrounds.

\item \textbf{Selection-principle frameworks.}
  The Penrose--Weyl curvature hypothesis (Penrose, \emph{Cycles of Time},
  2010) and the Past Hypothesis (Albert 2000) are the standard frameworks
  for the cosmological arrow of time.  Theorem T3 of ECI gives an algebraic
  correlate, but the low-complexity boundary condition remains a posit in
  all frameworks, including ECI.

\item \textbf{CLPW observer algebras.}
  Chandrasekaran--Penington--Witten 2023 (arXiv:2206.10780) is the parent
  framework for the type-II$_\infty$ construction.  Witten 2022
  (arXiv:2112.12828) is the foundational paper on the crossed-product
  algebra in de Sitter.  These references establish what \emph{can} be
  proven at the observer-algebra level; ECI extends them to FRW backgrounds
  but does not supersede their scope.

\end{itemize}

% ---------------------------------------------------------------------------
% Local bibliography entries referenced above.
% Add these to eciNotes.bib / eci.bib if not already present.
% ---------------------------------------------------------------------------
% @misc{ECI_cosmo_const_audit,
%   author = {Remondi\`ere, Kevin},
%   title  = {Honest reality check --- cosmological constant and arrow-of-time
%             claims},
%   year   = {2026},
%   note   = {Internal audit note,
%             \texttt{notes/cosmo\_const\_aot\_honest\_2026\_05\_02.md}}
% }
